\ifdefined\gkver\else\def\gkver{student}\fi
\documentclass[\gkver]{gaokaozhenti}
\begin{document}

\chapter{2012年安徽理综}
%% paperType: 理综物理部分

\begin{enumerate}
\item
%% number: 14
%% paperName: 2012年安徽理综第 14 题 6 分
%% typeId: 1
%% type: 单选题
%% chapter: 曲线运动
%% point: 人造地球卫星
%% method: 无
%% score: 6
%% degree: 650
%% duplicateId: 0
%% body:
我国发射的“天宫一号”和“神州八号”在对接前，“天宫一号”的运行轨道高度为$350\Ukm$，“神州八号”的运行轨道高度为$343\Ukm$。它们的运行轨道均视为圆周，则\xzanswer{B}
\fourchoices[answer=B]
{“天宫一号”比“神州八号”速度大}
{“天宫一号”比“神州八号”周期大}
{“天宫一号”比“神州八号”角速度大}
{“天宫一号”比“神州八号”加速度大}
%% answer: B
%% memo:
\memoanswer{
由万有引力提供航天器做圆周运动的向心力得 $G\frac{Mm}{r^{2}}=m\frac{v^{2}}{r}=m\omega^{2}r=m\frac{4\pi^{2}}{T^{2}}r=ma_{n}$，所以 $v=\sqrt{\frac{GM}{r}}$、$T=\sqrt{\frac{4\pi^{2}r^{3}}{GM}}$、$\omega=\sqrt{\frac{GM}{r^{3}}}$、$a_{n}=\frac{GM}{r^{2}}$。而“天宫一号”轨道半径 $r_{\text{天}}$ 比“神州八号”轨道半径 $r_{\text{神}}$ 大，故周期大。

正确选项：B。
}

\item
%% number: 15
%% paperName: 2012年安徽理综第 15 题 6 分
%% typeId: 1
%% type: 单选题
%% chapter: 机械波
%% point: 波的图象
%% method: 无
%% score: 6
%% degree: 650
%% duplicateId: 0
%% body:
一列简谐横波沿$x$轴正方向传播，在$t=0\Us$时波形如图1所示，已知波速为$10\Ums$，则$t=0.1\Us$时正确的波形是图2中的\xzanswer{C}
\onepicture[h=2.6cm, align=right]{figs/15.png}
\fourchoices[answer=C, ispicture=true, h=2.2cm]
{figs/15a.png}
{figs/15b.png}
{figs/15c.png}
{figs/15d.png}
%% answer: C
%% memo:
\memoanswer{
由图1可得波长 $\lambda=4.0\Um$，其周期 $T=\lambda/v=0.4\Us$。而 $t=0.1\Us=T/4$，波沿 $x$ 轴正方向传播，即图1的波形向 $x$ 轴正方向移动 $1/4$ 波长，得到图2的 C 图。

正确选项：C。
}

\item
%% number: 16
%% paperName: 2012年安徽理综第 16 题 6 分
%% typeId: 1
%% type: 单选题
%% chapter: 曲线运动
%% point: 竖直平面圆周运动
%% method: 无
%% score: 6
%% degree: 650
%% duplicateId: 0
%% body:
如图所示，在竖直平面内有一个半径为$R$的圆弧轨道。半径OA水平、OB竖直，一个质量为$m$的小球自A正上方P点由静止开始自由下落，小球沿轨道到达最高点B时恰好对轨道没有压力，已知$AP=2R$，重力加速度为$g$，则小球从P到B的运动过程中\xzanswer{D}
\onepicture[w=4.0cm]{\includesvg{figs/16}}
\fourchoices[answer=D]
{重力做功$2mgR$}
{机械能减少$mgR$}
{合外力做功$mgR$}
{克服摩擦力做功$\frac{1}{2}mgR$}
%% answer: D
%% memo:
\memoanswer{
小球沿轨道到达最高点 B 时恰好对轨道没有压力，由牛顿第二定律得 $mg=m\frac{v_{B}^{2}}{R}$；小球从 P 到 B 的运动过程中，由动能定理得 $W_{G}+W_{f}=\frac{1}{2}mv_{B}^{2}$。重力做功 $W_{G}=mgR$，合外力做功 $W=\frac{1}{2}mv_{B}^{2}=\frac{1}{2}mgR$，摩擦力做的功为 $W_{f}=-\frac{1}{2}mgR$，即克服摩擦力做功 $\frac{1}{2}mgR$，机械能减少 $\frac{1}{2}mgR$。

正确选项：D。
}

\item
%% number: 17
%% paperName: 2012年安徽理综第 17 题 6 分
%% typeId: 1
%% type: 单选题
%% chapter: 运动和力
%% point: 牛顿第二定律
%% method: 无
%% score: 6
%% degree: 450
%% duplicateId: 0
%% body:
如图所示，放在固定斜面上的物块以加速度$a$沿斜面匀加速下滑，若在物块上再施加一个竖直向下的恒力$F$，则\xzanswer{C}
\onepicture{figs/17.png}
\fourchoices[answer=C]
{物块可能匀速下滑}
{物块将以加速度$a$匀加速下滑}
{物块将以大于$a$的加速度匀加速下滑}
{物块将以小于$a$的加速度匀加速下滑}
%% answer: C
%% memo:
\memoanswer{
设斜面的倾角为 $\alpha$，物块与斜面间动摩擦因数为 $\mu$，施加一个竖直向下的恒力 $F$ 时加速度为 $a'$。

根据牛顿第二定律，不施加恒力 $F$ 时：$mg\sin\alpha-\mu mg\cos\alpha=ma$，得 $a=g(\sin\alpha-\mu\cos\alpha)$；施加一个竖直向下的恒力 $F$ 时：$(mg+F)\sin\alpha-\mu(mg+F)\cos\alpha=ma'$，得 $a'=\left(g+\frac{F}{m}\right)(\sin\alpha-\mu\cos\alpha)$。

比较可知 $a'>a$，故正确选项：C。
}

\item
%% number: 18
%% paperName: 2012年安徽理综第 18 题 6 分
%% typeId: 1
%% type: 单选题
%% chapter: 电场
%% point: 电势差与电场强度的关系
%% method: 无
%% score: 6
%% degree: 450
%% duplicateId: 0
%% body:
如图所示，在平面直角坐标系中，有方向平行于坐标平面的匀强电场，其中坐标原点$O$处的电势为$0\UV$，点A处的电势为$6\UV$，点B处的电势为$3\UV$，则电场强度的大小为\xzanswer{A}
\onepicture[w=6.5cm]{\includesvg{figs/18}}
\fourchoices[answer=A]
{$200\UVm$}
{$200\sqrt{3}\UVm$}
{$100\UVm$}
{$100\sqrt{3}\UVm$}
%% answer: A
%% memo:
\memoanswer{
由于 $\varphi_{A}=6\UV$、$\varphi_{B}=3\UV$、$\varphi_{O}=0\UV$，且是匀强电场，则在 $(3,0)$ 的点的电势为 $3\UV$。匀强电场的方向垂直于点 $(3,0)$ 与 B 点的连线，由几何关系可得 O 到点 $(3,0)$ 与 B 点的连线的距离 $d=1.5\Ucm$，匀强电场的电场强度 $E=U_{BO}/d=200\UVm$。

正确选项：A。
}

\item
%% number: 19
%% paperName: 2012年安徽理综第 19 题 6 分
%% typeId: 1
%% type: 单选题
%% chapter: 磁场
%% point: 带电粒子在磁场中的运动
%% method: 无
%% score: 6
%% degree: 450
%% duplicateId: 0
%% body:
如图所示，圆形区域内有垂直于纸面向里的匀强磁场，一个带电粒子以速度$v$从A点沿直径AOB方向射入磁场，经过$\Delta t$时间从C点射出磁场，OC与OB成$60{^\circ}$角。现将带电粒子的速度变为$v/3$，仍从A点射入磁场，不计重力，则粒子在磁场中的运动时间变为\xzanswer{B}
\onepicture{figs/19.png}
\fourchoices[answer=B]
{$\frac{1}{2}\Delta t$}
{$2\Delta t$}
{$\frac{1}{3}\Delta t$}
{$3\Delta t$}
%% answer: B
%% memo:
\memoanswer{
由牛顿第二定律 $qvB=m\frac{v^{2}}{r}$ 及匀速圆周运动 $T=\frac{2\pi r}{v}$ 得 $r=\frac{mv}{qB}$，$T=\frac{2\pi m}{qB}$。

由图可得以速度 $v$ 从 A 点沿直径 AOB 方向射入磁场，经过 $\Delta t=T/6$ 从 C 点射出磁场，轨道半径 $r=\sqrt{3}\,AO$；速度变为 $v/3$ 时，运动半径是 $r/3=\frac{\sqrt{3}}{3}AO$，由几何关系可得在磁场中运动转过的圆心角为 $120{^\circ}$，运动时间为 $T/3$，即 $2\Delta t$。

正确选项：B。
}

\item
%% number: 20
%% paperName: 2012年安徽理综第 20 题 6 分
%% typeId: 1
%% type: 单选题
%% chapter: 电场
%% point: 电场叠加
%% method: 极限
%% score: 6
%% degree: 450
%% duplicateId: 0
%% body:
如图1所示，半径为$R$的均匀带电圆形平板，单位面积带电量为$\sigma$，其轴线上任意一点P（坐标为$x$）的电场强度可以由库仑定律和电场强度的叠加原理求出：$E=2\pi k\sigma\left[1-\frac{x}{(R^{2}+x^{2})^{1/2}}\right]$，方向沿$x$轴。现考虑单位面积带电量为$\sigma_{0}$的无限大均匀带电平板，从其中间挖去一半径为$r$的圆形板，如图2所示。则圆孔轴线上任意一点Q（坐标为$x$）的电场强度为\xzanswer{A}
\onepicture{figs/20.png}
\fourchoices[answer=A]
{$2\pi k\sigma_{0}\frac{x}{(r^{2}+x^{2})^{1/2}}$}
{$2\pi k\sigma_{0}\frac{r}{(r^{2}+x^{2})^{1/2}}$}
{$2\pi k\sigma_{0}\frac{x}{r}$}
{$2\pi k\sigma_{0}\frac{r}{x}$}
%% answer: A
%% memo:
\memoanswer{
无限大均匀带电平板周围的电场是垂直于平板的匀强电场，即电场强度处处相等等于圆板中心轴上 $x=0$ 时的电场强度，由题中信息可得单位面积带电量为 $\sigma_{0}$ 的无限大均匀带电平板场强为 $E=2\pi k\sigma_{0}$。而半径为 $r$ 的圆板在其轴线上 Q 点的电场场强为 $E'=2\pi k\sigma_{0}\left[1-\frac{x}{(r^{2}+x^{2})^{1/2}}\right]$，由电场叠加原理可得图2中 Q（坐标为 $x$）的电场强度为 $E$ 和 $E'$ 的矢量和，即 $E-E'=2\pi k\sigma_{0}\frac{x}{(r^{2}+x^{2})^{1/2}}$。

正确选项：A。
}

\item
%% number: 21
%% paperName: 2012年安徽理综第 21 题 10 分
%% typeId: 4
%% type: 实验题
%% chapter: 运动和力
%% point: 验证牛顿第二定律
%% method: 无
%% score: 10
%% degree: 450
%% duplicateId: 0
%% body:
图1为“验证牛顿第二定律”的实验装置示意图。砂和砂桶的总质量为$m$，小车和砝码的总质量为$M$。实验中用砂和砂桶总重力的大小作为细线对小车拉力的大小。
\begin{enumerate}
	\item
	实验中，为了使细线对小车的拉力等于小车所受的合外力，先调节长木板一端滑轮的高度，使细线与长木板平行。接下来还需要进行的一项操作是（\quad）

	（A）将长木板水平放置，让小车连着已经穿过打点计时器的纸带，给打点计时器通电，调节$m$的大小，使小车在砂和砂桶的牵引下运动，从打出的纸带判断小车是否做匀速运动。

	（B）将长木板的一端垫起适当的高度，让小车连着已经穿过打点计时器的纸带，撤去砂和砂桶，给打点计时器通电，轻推小车，从打出的纸带判断小车是否做匀速运动。

	（C）将长木板的一端垫起适当的高度，撤去纸带以及砂和砂桶，轻推小车，观察判断小车是否做匀速运动。

	\item
	实验中要进行质量$m$和$M$的选取，以下最合理的一组是（\quad）

	（A）$M=20\Ug$，$m=10\Ug$、$15\Ug$、$20\Ug$、$25\Ug$、$30\Ug$、$40\Ug$

	（B）$M=200\Ug$，$m=20\Ug$、$40\Ug$、$60\Ug$、$80\Ug$、$100\Ug$、$120\Ug$

	（C）$M=400\Ug$，$m=10\Ug$、$15\Ug$、$20\Ug$、$25\Ug$、$30\Ug$、$40\Ug$

	（D）$M=400\Ug$，$m=20\Ug$、$40\Ug$、$60\Ug$、$80\Ug$、$100\Ug$、$120\Ug$

	\item
	图2是实验中得到的一条纸带，A、B、C、D、E、F、G为7个相邻的计数点，相邻的两个计数点之间还有四个点未画出。量出相邻的计数点之间的距离分别为$s_{AB}=4.22\Ucm$、$s_{BC}=4.65\Ucm$、$s_{CD}=5.08\Ucm$、$s_{DE}=5.49\Ucm$、$s_{EF}=5.91\Ucm$、$s_{FG}=6.34\Ucm$。已知打点计时器的工作频率为$50\UHz$，则小车的加速度$a=$\_\_\_\_\_\_\_\_\_$\Umsq$（结果保留2位有效数字）。
\end{enumerate}
\twopicture[labelA=2012安徽理综21a, labelB=2012安徽理综21b, showlabel=false, align=right]
{figs/21a.png}
{figs/21b.png}
%% answer: 
\jdanswer{
\begin{enumerate}
	\item
	B
	\item
	C
	\item
	0.42
\end{enumerate}
}
%% memo:
\memoanswer{
（1）平衡摩擦阻力是通过 B 的方法来实现的，故选 B。

（2）由于本实验中要求砂和砂桶的质量 $m$ 远小于小车和砝码的质量 $M$，故选 C 合理。

（3）由于连续相等时间内的位移差 $\Delta s=0.42\Ucm$，由匀变速运动规律 $\Delta s=aT^{2}$，且 $T=5\times0.02\Us=0.1\Us$，所以 $a=\Delta s/T^{2}=0.42\Umsq$。
}

\item
%% number: 21
%% paperName: 2012年安徽理综第 21 题 8 分
%% typeId: 4
%% type: 实验题
%% chapter: 电路
%% point: 伏安法测电阻
%% method: 无
%% score: 8
%% degree: 650
%% duplicateId: 0
%% body:
图为“测绘小灯泡伏安特性曲线”实验的实物电路图，已知小灯泡额定电压为$2.5\UV$。
\begin{enumerate}
	\item
	完成下列实验步骤：

	①闭合开关前，调节滑动变阻器的滑片，\_\_\_\_\_\_\_\_\_\_\_\_\_\_；

	②闭合开关后，逐渐移动变阻器的滑片，\_\_\_\_\_\_\_\_\_\_\_\_\_\_；

	③断开开关，$\ldots\ldots$。根据实验数据在方格纸上作出小灯泡灯丝的伏安特性曲线。
	\item
	在虚线框中画出与实物电路相应的电路图。
\end{enumerate}
\onepicture[align=right]{figs/21.png}
%% answer: 
\jdanswer{
\begin{enumerate}
	\item
	①使滑动变阻器的滑片靠近变阻器左端的接线柱（使小灯泡两端电压为 0）；②增加小灯泡两端的电压，记录电流表和电压表的多组读数，直至电压达到额定电压。
	\item
	如图所示。
\end{enumerate}
}
%% memo:
\memoanswer{
\onepicture[h=3cm, align=right]{figs/21_an.png}

（1）①由于滑动变阻器的分压作用，使开始实验时小灯泡两端电压为 0，应使滑动变阻器的滑片靠近变阻器左端的接线柱。②描绘小灯泡灯丝的伏安特性曲线，必须测量多组数据，即增加小灯泡两端的电压，记录电流表和电压表的多组读数，直至电压达到额定电压。

（2）与实物电路相应的电路图如图所示。
}

\item
%% number: 22
%% paperName: 2012年安徽理综第 22 题 14 分
%% typeId: 6
%% type: 计算题
%% chapter: 运动和力
%% point: 牛顿定律的应用
%% method: 无
%% score: 14
%% degree: 650
%% duplicateId: 0
%% body:
质量为$0.1\Ukg$的弹性球从空中某高度由静止开始下落，该下落过程对应的$v-t$图象如图所示。球与水平地面相碰后离开地面时的速度大小为碰撞前的$3/4$。设球受到的空气阻力大小恒为$f$，取$g=10\Umsq$，求：
\begin{enumerate}
	\item
	弹性球受到的空气阻力$f$的大小；
	\item
	弹性球第一次碰撞后反弹的高度$h$。
\end{enumerate}
\onepicture[align=right]{figs/22.png}
%% answer: 
\jdanswer{
\begin{enumerate}
	\item
	$0.2\UN$
	\item
	$\frac{3}{8}\Um$
\end{enumerate}
}
%% memo:
\memoanswer{
（1）设弹性球第一次下落过程中的加速度大小为 $a_{1}$，由图知 $a_{1}=\frac{\Delta v}{\Delta t}=\frac{4}{0.5}\Umsq=8\Umsq$。

根据牛顿第二定律，得 $mg-f=ma_{1}$，解得 $f=m(g-a_{1})=0.2\UN$。

（2）由图知弹性球第一次到达地面时的速度大小为 $v_{1}=4\Ums$，设球第一次离开地面时的速度为 $v_{2}$，则 $v_{2}=\frac{3}{4}v_{1}=3\Ums$。第一次离开地面后，设上升过程中球的加速度大小为 $a_{2}$，则 $mg+f=ma_{2}$，得 $a_{2}=12\Umsq$。于是有 $0-v_{2}^{2}=-2a_{2}h$，解得 $h=\frac{3}{8}\Um$。
}

\item
%% number: 23
%% paperName: 2012年安徽理综第 23 题 16 分
%% typeId: 6
%% type: 计算题
%% chapter: 交流电
%% point: 交流电
%% method: 无
%% score: 16
%% degree: 450
%% duplicateId: 0
%% body:
图1是交流发电机模型示意图。在磁感应强度为$B$的匀强磁场中，有一矩形线圈abcd可绕线圈平面内垂直于磁感线的$OO'$轴转动，由线圈引出的导线ae和df分别与两个跟线圈一起绕$OO'$转动的金属圆环相连接，金属圆环又分别与两个固定的电刷保持滑动接触，这样矩形线圈在转动中就可以保持和外电路电阻$R$形成闭合电路。图2是线圈的主视图，导线ab和cd分别用它们的横截面来表示。已知ab长度为$L_{1}$，bc长度为$L_{2}$，线圈以恒定角速度$\omega$逆时针转动。（只考虑单匝线圈）
\begin{enumerate}
	\item
	线圈平面处于中性面位置时开始计时，试推导$t$时刻整个线圈中的感应电动势$e_{1}$的表达式；
	\item
	线圈平面处于与中性面成$\varphi_{0}$夹角位置时开始计时，如图3所示，试写出$t$时刻整个线圈中的感应电动势$e_{2}$的表达式；
	\item
	若线圈电阻为$r$，求线圈每转动一周电阻$R$上产生的焦耳热。（其它电阻均不计）
\end{enumerate}
\onepicture[align=right]{figs/23.png}
%% answer: 
\jdanswer{
\begin{enumerate}
	\item
	$e_{1}=BL_{1}L_{2}\omega\sin\omega t$
	\item
	$e_{2}=BL_{1}L_{2}\omega\sin(\omega t+\varphi_{0})$
	\item
	$Q_{R}=\pi R\omega\left(\frac{BL_{1}L_{2}}{R+r}\right)^{2}$
\end{enumerate}
}
%% memo:
\memoanswer{
（1）矩形线圈 abcd 转动过程中，只有 ab 和 cd 切割磁感线，设 ab 和 cd 的转动速度为 $v$，则 $v=\omega\frac{L_{2}}{2}$。在 $t$ 时刻，导线 ab 和 cd 因切割磁感线而产生的感应电动势均为 $E_{1}=BL_{1}v$。由图可知 $v_{\perp}=v\sin\omega t$。则整个线圈的感应电动势为 $e_{1}=2E_{1}=BL_{1}L_{2}\omega\sin\omega t$。

（2）当线圈由图3位置开始运动时，在 $t$ 时刻整个线圈的感应电动势为 $e_{2}=BL_{1}L_{2}\omega\sin(\omega t+\varphi_{0})$。

（3）由闭合电路欧姆定律可知 $I=\frac{E}{R+r}$，这里的 $E$ 为线圈产生的电动势的有效值，$E=\frac{E_{m}}{\sqrt{2}}=\frac{BL_{1}L_{2}\omega}{\sqrt{2}}$。则线圈转动一周在 $R$ 上产生的焦耳热为 $Q_{R}=I^{2}RT$，其中 $T=\frac{2\pi}{\omega}$。于是 $Q_{R}=\pi R\omega\left(\frac{BL_{1}L_{2}}{R+r}\right)^{2}$。
}

\item
%% number: 24
%% paperName: 2012年安徽理综第 24 题 20 分
%% typeId: 6
%% type: 计算题
%% chapter: 运动和力
%% point: 动量守恒定律
%% method: 无
%% score: 20
%% degree: 250
%% duplicateId: 0
%% body:
如图所示，装置的左边是足够长的光滑水平面，一轻质弹簧左端固定，右端连接着质量$M=2\Ukg$的小物块A。装置的中间是水平传送带，它与左右两边的台面等高，并能平滑对接。传送带始终以$u=2\Ums$的速率逆时针转动。装置的右边是一光滑的曲面，质量$m=1\Ukg$的小物块B从其上距水平台面$h=1.0\Um$处由静止释放。已知物块B与传送带之间的摩擦因数$\mu=0.2$，$l=1.0\Um$。设物块A、B中间发生的是对心弹性碰撞，第一次碰撞前物块A静止且处于平衡状态。取$g=10\Umsq$。
\begin{enumerate}
	\item
	求物块B与物块A第一次碰撞前速度大小；
	\item
	通过计算说明物块B与物块A第一次碰撞后能否运动到右边曲面上？
	\item
	如果物块A、B每次碰撞后，物块A再回到平衡位置时都会立即被锁定，而当他们再次碰撞前锁定被解除，试求出物块B第$n$次碰撞后的运动速度大小。
\end{enumerate}
\onepicture[align=right]{figs/24.png}
%% answer: 
\jdanswer{
\begin{enumerate}
	\item
	$4\Ums$
	\item
	不能
	\item
	$\left(\frac{1}{3}\right)^{n}v$
\end{enumerate}
}
%% memo:
\memoanswer{
（1）设物块 B 沿光滑曲面下滑到水平位置时的速度大小为 $v_{0}$，由机械能守恒知 $mgh=\frac{1}{2}mv_{0}^{2}$，得 $v_{0}=\sqrt{2gh}$。设物块 B 在传送带上滑动过程中因受摩擦力所产生的加速度大小为 $a$，则 $\mu mg=ma$。设物块 B 通过传送带后运动速度大小为 $v$，有 $v^{2}-v_{0}^{2}=-2al$。结合以上各式解得 $v=4\Ums$。由于 $v>u=2\Ums$，所以 $v=4\Ums$ 即为物块 B 与物块 A 第一次碰撞前的速度大小。

（2）设物块 A、B 第一次碰撞后的速度分别为 $V$、$v_{1}$，取向右为正方向，由弹性碰撞知 $-mv=mv_{1}+MV$，$\frac{1}{2}mv^{2}=\frac{1}{2}mv_{1}^{2}+\frac{1}{2}MV^{2}$，解得 $v_{1}=\frac{1}{3}v=\frac{4}{3}\Ums$。即碰撞后物块 B 在水平台面向右匀速运动。

设物块 B 在传送带上向右运动的最大位移为 $l'$，则 $0-v_{1}^{2}=-2al'$，得 $l'=\frac{4}{9}\Um<1\Um$。所以物块 B 不能通过传送带运动到右边的曲面上。

（3）当物块 B 在传送带上向右运动的速度为零时，将会沿传送带向左加速。可以判断，物块 B 运动到左边台面时的速度大小为 $v_{1}$，继而与物块 A 发生第二次碰撞。设第二次碰撞后物块 B 速度大小为 $v_{2}$，同上计算可知 $v_{2}=\frac{1}{3}v_{1}=\left(\frac{1}{3}\right)^{2}v$。物块 B 与物块 A 第三次碰撞、第四次碰撞……碰撞后物块 B 的速度大小依次为 $v_{3}=\frac{1}{3}v_{2}=\left(\frac{1}{3}\right)^{3}v$，$v_{4}=\frac{1}{3}v_{3}=\left(\frac{1}{3}\right)^{4}v$，$\cdots\cdots$。则第 $n$ 次碰撞后物块 B 的速度大小为 $v_{n}=\left(\frac{1}{3}\right)^{n}v$。
}

\end{enumerate}

\end{document}
